% front/08_eng_abstract.tex

\clearpage
\newcommand{\EnglishAbstractTitle}{\mbox{영문초록}}
\newcommand{\EnglishAbstractTOC}{\PlainTOCtext{영문초록}}
\let\EngAbstractHeading\EnglishAbstractTitle
\phantomsection
\addcontentsline{toc}{chapter}{\EnglishAbstractTOC}

% ① "Abstract" title (left aligned, 16pt bold)
{\fontsize{16}{20}\selectfont\bfseries\EnglishAbstractTitle}\par
\vspace{15mm}  % 1.5cm

\begin{center}
  \begin{minipage}{0.9\textwidth}
    \centering
    {\fontsize{16}{20}\selectfont\bfseries
      A Dynamic RTO Framework\par
      Based on Observability and Uncertainty\par}
  \end{minipage}
\end{center}

% ③ Name/department block – right aligned, inset by 2cm
\vspace{15mm}  % 1.5cm

\begin{flushright}
  \hspace*{-2cm}
  {\fontsize{14}{22}\selectfont
    \setlength{\baselineskip}{22pt}
    LEE,\ CHANGHO\par
    Department of Disaster and Safety Management\par
    Graduate School of Soongsil University\par
  }
\end{flushright}

\vspace{15mm}  % 1.5cm

% ④ Main text (11pt, line spacing 200%)
{\fontsize{11}{22}\selectfont
\setstretch{2.0}\justifying
\setlength{\parindent}{2em}

In business continuity management (BCM), the Recovery Time Objective (RTO) is determined as a fixed value during business impact analysis (BIA) and remains unchanged regardless of subsequent changes in the operational environment. This static RTO approach has structural limitations in that it fails to account for dynamic fluctuations in system observability, the probabilistic nature of uncertainty inherent in recovery processes, and heavy-tail risks from extreme events such as cyberattacks and compound disasters.

To address these limitations, this study proposes a Dynamic RTO (DRTO) framework comprising seven components: (i) a mathematical model that formulates observability and uncertainty via individual sigmoid bounding functions, (ii) four operational conditions (C0--C3) spanning from a static baseline to heavy-tail uncertainty, (iii) Monte Carlo simulation with $n = 10{,}000$ per condition, (iv) multi-dimensional comparative analysis including effect sizes, CDF-crossover-based split regression, and CVaR tail-risk quantification, (v) robustness verification through multi-seed testing ($10^8$ runs), (vi) a six-stage multi-layer validation framework (L0--L5) combining Golden Compare foundation verification, ten hypothesis tests, and expert evaluation, and (vii) generalizability verification across six industry scenarios. The core equation is defined as $\DRTO = R_0 + E_{O,\mathrm{bnd}} + E_{U,\mathrm{bnd}}$, where the observability effect $E_{O,\mathrm{bnd}} = -R_0 \cdot S(z_O)$ adjusts recovery time downward from the baseline $R_0$, and the uncertainty effect $E_{U,\mathrm{bnd}} = (R_{\max} - R_0) \cdot S(z_U)$ adjusts it upward. The modified sigmoid $S(z) = 2/(1 + e^{-z}) - 1$ asymptotically guarantees that $\DRTO$ remains within the physical bounds $(0,\; R_{\max})$. The default parameters are $R_0 = 6.0$ hours, $R_{\max} = 24.0$ hours ($= 4R_0$), and $\kappa = 1.2$. \fpara{The value $\kappa = 1.2$ was derived from grid search in environment C1 (observability-only, Gaussian baseline) and is adopted not as an \emph{environment-specific optimum} but as a \emph{provisional and conservative operational default} for general BCM operations. Its applicability in environments C2/C3 is interpreted as \emph{suggesting operational adequacy} through the integrated PASS of hypotheses H2--H10; the \emph{investigation of the existence} of an environment-invariant $\kappa$ is deferred to future research.}

The model defines four operational conditions in terms of the efficiency ratio $\eta = \DRTO / R_0$. \chg{The static baseline C0 has $\eta = 1.000$, the observability-only C1 has $\eta \leq 1.000$, the Gaussian-uncertainty C2 has $\eta \geq \eta_{\mathrm{C1}}$, and the Pareto-uncertainty C3 has $\eta \geq \eta_{\mathrm{C2}}$ in the tail region.} Both distributions share a unified expected value of $3.0$, enabling fair comparison of distributional shape effects.

Monte Carlo simulation with $n = 10{,}000$ per condition revealed a two-stage branching structure in the efficiency ratio $\eta = \DRTO / R_0$: the common path C0 ($1.000$) $\to$ C1 ($0.853$) showed a mean reduction of $14.7\%$ (Cohen's $d = -1.180$), confirming that observability significantly decreases recovery time. From C1, two parallel branches emerged: C2 ($1.682$, Gaussian) and C3 ($1.514$, Pareto), both elevated by uncertainty premiums. Despite the common mean design, C2 and C3 exhibited marked differences at extreme percentiles: at the 99th percentile, C3 $\DRTO$ reached approximately $13.2$ hours compared to C2 at approximately $8.4$ hours---a $57\%$ gap attributable to the Pareto distribution's $U_{\mathrm{raw}}$ 99th percentile being approximately $4.1$ times larger than the Gaussian ($21.85$ vs.\ $5.33$).

Using the CDF crossover point at P85.8 ($\eta \approx 2.42$), split regression between Core ($\leq$ P85.8) and Tail ($>$ P85.8) regions revealed a dual-channel attenuation effect: the regression coefficient of observability weight $k_O$ attenuated from $-0.798$ (Core) to $-0.415$ (Tail), a $0.52\times$ reduction indicating that observability's recovery-shortening effect is structurally weakened under extreme uncertainty. CVaR$_{99}$ for C3 exceeded that of C2 by $24.0\%$, quantitatively demonstrating that heavy-tail risk disproportionately amplifies recovery burden in extreme regions.

The framework was validated through a six-stage multi-layer verification system (L0--L5). Foundation verification (L0, Golden Compare) confirmed implementation correctness via 47 deterministic checkpoints (100\% PASS). All ten hypotheses (H1--H10) were accepted. Expert evaluation by five BCM/DR specialists (mean $15.2$ years of experience) yielded a quantitative survey mean of $4.46/5.0$ ($p = 0.011$, Cohen's $d = 1.65$) with high internal consistency (Cronbach's $\alpha = 0.89$, S-CVI/Ave$\,= 0.985$), and in-depth interviews confirmed $\geq 80\%$ agreement on all 15 sub-themes across 5 major themes ($100\%$ met). Full quadratic regression (including interaction terms) achieved $R^2 = 1.000$ for C1, $R^2 = 0.998$ for C2, and $R^2 = 0.858$ for C3 overall, with the P85.8 split yielding Core $R^2 = 0.999$ and Tail $R^2 = 0.710$, validating the effectiveness of split regression. Additionally, six industry scenarios ($R_0$: $0.5$--$8.0$ hours) showed a consistent P99 ratio averaging $1.57\times$, confirming the model's general applicability.

This research contributes to business continuity management by integrating observability benefits and uncertainty penalties through individual bounding, and by elucidating the disproportionate impact of heavy-tail distributions and the dual-channel attenuation mechanism. It thereby provides a statistically grounded methodology for dynamic RTO adjustment, a regression-based DRTO computation tool, and---grounded in the dual-channel attenuation finding---resource-allocation principles and stress-test design criteria. \fpara{As an acknowledged academic limitation, the environment-invariant universal optimality of $\kappa = 1.2$ is not proven in this study and is deferred to future research.}

\vspace{10mm}
\noindent\textbf{Keywords:} Recovery Time Objective(RTO), Dynamic RTO(DRTO), Business Continuity Management(BCM), Observability, Uncertainty Quantification, Sigmoid Bounding, Heavy-tailed Distribution

}
